\documentclass[12pt,a4paper]{article}

% Encoding and fonts
\usepackage[utf8]{inputenc}
\usepackage[T1]{fontenc}
\usepackage{lmodern}
\usepackage{textcomp}

% Page layout
\usepackage[top=1in, bottom=1in, left=1in, right=1in]{geometry}

% Typography
\usepackage{microtype}
\usepackage{parskip}

% Language
\usepackage[english]{babel}

% Headers and footers
\usepackage{fancyhdr}
\pagestyle{fancy}
\fancyhf{}
\fancyhead[L]{\leftmark}
\fancyhead[R]{\thepage}
\fancyfoot[C]{\thepage}
\renewcommand{\headrulewidth}{0.4pt}
\renewcommand{\footrulewidth}{0pt}

% Section styling
\usepackage{titlesec}
\titleformat{\section}{\Large\bfseries}{\thesection}{1em}{}
\titleformat{\subsection}{\large\bfseries}{\thesubsection}{1em}{}
\titleformat{\subsubsection}{\normalsize\bfseries}{\thesubsubsection}{1em}{}
\titlespacing*{\section}{0pt}{2.5ex plus 1ex minus .2ex}{1.5ex plus .2ex}
\titlespacing*{\subsection}{0pt}{2ex plus 1ex minus .2ex}{1ex plus .2ex}
\titlespacing*{\subsubsection}{0pt}{1.5ex plus 1ex minus .2ex}{0.8ex plus .2ex}

% List formatting
\usepackage{enumitem}
\setlist[itemize]{topsep=0.5ex, itemsep=0.25ex, parsep=0pt}
\setlist[enumerate]{topsep=0.5ex, itemsep=0.25ex, parsep=0pt}

% Essential packages
\usepackage{graphicx}
\usepackage[table]{xcolor}
\usepackage{booktabs}
\usepackage{array}
\usepackage{tabularx}
\usepackage{longtable}
\usepackage{amsmath}
\usepackage{amssymb}
\usepackage[normalem]{ulem}
\rowcolors{2}{gray!10}{white}
\renewcommand{\arraystretch}{1.3}

% Cross-references
\usepackage{hyperref}
\usepackage{cleveref}

% Custom commands
\newcommand{\pilcrow}{\S}

% Hyperref setup
\hypersetup{
    colorlinks=true,
    linkcolor=blue,
    filecolor=magenta,
    urlcolor=cyan,
}

% Code listings
\usepackage{listings}
\definecolor{codebg}{RGB}{248,248,248}
\definecolor{codeframe}{RGB}{200,200,200}
\definecolor{codekeyword}{RGB}{0,0,180}
\definecolor{codecomment}{RGB}{0,128,0}
\definecolor{codestring}{RGB}{163,21,21}
\lstset{
    basicstyle=\ttfamily\small,
    breaklines=true,
    breakatwhitespace=false,
    frame=single,
    framerule=0.5pt,
    rulecolor=\color{codeframe},
    backgroundcolor=\color{codebg},
    keepspaces=true,
    showstringspaces=false,
    tabsize=4,
    xleftmargin=1em,
    xrightmargin=1em,
    aboveskip=1em,
    belowskip=1em,
    keywordstyle=\color{codekeyword}\bfseries,
    commentstyle=\color{codecomment}\itshape,
    stringstyle=\color{codestring},
    literate={_}{{\textunderscore}}1
             {-}{{\textminus}}1
             {<}{{\textless}}1
             {>}{{\textgreater}}1
             {~}{{\textasciitilde}}1
             {^}{{\textasciicircum}}1
}

% YAML language definition for listings
\lstdefinelanguage{YAML}{
    keywords={true,false,null,y,n},
    keywordstyle=\color{blue}\bfseries,
    sensitive=false,
    comment=[l]{\#},
    morecomment=[s]{/*}{*/},
    commentstyle=\color{gray}\ttfamily,
    stringstyle=\color{red}\ttfamily,
    morestring=[b]'',
    morestring=[b]'',
}

% TOML language definition for listings
\lstdefinelanguage{TOML}{
    keywords={true,false},
    keywordstyle=\color{blue}\bfseries,
    sensitive=true,
    comment=[l]{\#},
    commentstyle=\color{gray}\ttfamily,
    stringstyle=\color{red}\ttfamily,
    morestring=[b]'',
    morestring=[b]'',
}

% Dockerfile language definition for listings
\lstdefinelanguage{Dockerfile}{
    keywords={FROM,RUN,CMD,LABEL,EXPOSE,ENV,ADD,COPY,ENTRYPOINT,VOLUME,USER,WORKDIR,ARG,ONBUILD,STOPSIGNAL,HEALTHCHECK,SHELL},
    keywordstyle=\color{blue}\bfseries,
    sensitive=false,
    comment=[l]{\#},
    commentstyle=\color{gray}\ttfamily,
    stringstyle=\color{red}\ttfamily,
    morestring=[b]'',
    morestring=[b]'',
}

% JSON language definition for listings
\lstdefinelanguage{JSON}{
    keywords={true,false,null},
    keywordstyle=\color{blue}\bfseries,
    sensitive=true,
    commentstyle=\color{gray}\ttfamily,
    stringstyle=\color{red}\ttfamily,
    morestring=[b]'',
}

% Code listing float environment
\usepackage{float}
\newfloat{listing}{htbp}{lol}[section]
\floatname{listing}{Listing}

% Admonition boxes
\usepackage{tcolorbox}
\tcbuselibrary{skins,breakable}

% Admonition style definitions
\newtcolorbox{admonitionbox}[2][]{
    enhanced,
    breakable,
    colback=#2!5,
    colframe=#2!80!black,
    fonttitle=\bfseries,
    title=#1,
    left=2mm,
    right=2mm,
}

% Synthon tuple boxes
\newtcolorbox{synthonbox}{
    enhanced,
    colback=blue!5,
    colframe=blue!75!black,
    fontupper=\large,
    center upper,
    boxrule=1pt,
    arc=4pt,
}

\begin{document}

\title{ALEPH: Architecture and Verification}
\date{\today}

\maketitle

\subsection{From Grammar to Kernel to Proof}

\begin{center}
\rule{0.5\textwidth}{0.4pt}
\end{center}

\subsection{The Imscribing Grammar}

The Imscribing Grammar (IG) is a twelve-primitive structural classification system. Every object --- mathematical, physical, computational, linguistic --- occupies a position in a twelve-dimensional type space whose coordinates are the primitive values. The twelve primitives are:

\begin{table}[htbp]
\centering
\small
\rowcolors{1}{white}{white}
\begin{tabularx}{\textwidth}{>{\centering\arraybackslash}X >{\centering\arraybackslash}X >{\centering\arraybackslash}X >{\centering\arraybackslash}X}
\toprule
\rowcolor{gray!25}\textbf{Index} & \textbf{Symbol} & \textbf{Name} & \textbf{Role} \\
\midrule
\rowcolors{1}{white}{gray!10}
0 & Ð & Dimensional depth & How many layers deep the object''s structure extends \\
1 & Þ & Topological type & The connectivity class of the object \\
2 & Ř & Relational mode & How the object relates to others of its kind \\
3 & $\Phi$ & Frobenius parity & The symmetry class: asym / sym / $\pm$ / \} (Frobenius-special) \\
4 & ƒ & Kinetic mode & The object''s dynamical regime \\
5 & Ç & Criticality kinetics & How criticality propagates through the object \\
6 & $\Gamma$ & Mediative structure & How the object participates in mediated relationships \\
7 & ɢ & Valency & The object''s coupling grain \\
8 & ⊙ & Self-modeling & The criticality / self-reference posture \\
9 & Ħ & Chirality & The handedness the object carries intrinsically \\
10 & $\Sigma$ & Symmetry type & Global symmetry class \\
11 & $\Omega$ & Winding & The topological winding invariant \\
\bottomrule
\end{tabularx}
\end{table}

The full type space has \textbf{17,280,000} positions (3³ $\times$ 4⁵ $\times$ 5⁴). Of these, 49 positions are occupied by structurally significant objects --- the catalog entries that appear in the IG symbol set.

\subsubsection{Tier Classification}

Five tiers partition the type space by structural depth:

\begin{table}[htbp]
\centering
\small
\rowcolors{1}{white}{white}
\begin{tabularx}{\textwidth}{>{\centering\arraybackslash}X >{\centering\arraybackslash}X >{\centering\arraybackslash}X}
\toprule
\rowcolor{gray!25}\textbf{Tier} & \textbf{Condition} & \textbf{Meaning} \\
\midrule
\rowcolors{1}{white}{gray!10}
O\_0 & ⊙ = sub or ⊙ = EP & No self-modeling, or exceptional-point collapse \\
O\_1 & ⊙ $\neq$ sub, $\Omega$ = 0 & Self-modeling but no winding \\
O\_2 & $\Omega$ $\neq$ 0, Ð $\in$ \{0,1,3\} & Wound, non-Frobenius \\
O\_2d & $\Omega$ $\neq$ 0, Ð = 2 & Wound, depth-2 variant \\
$O_{\infty}$ & ⊙ = c, $\Phi$ = \} & Critical self-modeling + Frobenius-special \\
\bottomrule
\end{tabularx}
\end{table}

$O_{\infty}$ is the fixed-point tier: an $O_{\infty}$ object is structurally self-consistent under its own operations. The defining algebraic property is the Frobenius condition:

\[\mu \circ \delta = \mathrm{id}\]

where $\delta$ is comultiplication (split) and $\mu$ is multiplication (fuse). An $O_{\infty}$ object can be split into two copies and fused back to itself without structural loss.

\begin{center}
\rule{0.5\textwidth}{0.4pt}
\end{center}

\subsection{The ALEPH Type System}

The ALEPH type system takes the twenty-two letters of the Hebrew alphabet as a natural, complete sample of the IG type space. Each letter is assigned a twelve-dimensional tuple, and its tier is derived from that tuple by the tier classification rules.

The twenty-two letters distribute across the tier hierarchy as follows:

\begin{table}[htbp]
\centering
\small
\rowcolors{1}{white}{white}
\begin{tabularx}{\textwidth}{>{\centering\arraybackslash}X >{\centering\arraybackslash}X >{\centering\arraybackslash}X}
\toprule
\rowcolor{gray!25}\textbf{Tier} & \textbf{Count} & \textbf{Letters} \\
\midrule
\rowcolors{1}{white}{gray!10}
O\_\textbackslash{}infty & 3 & ו (vav), מ  (mem), ש  (shin) \\
O\_2\_\{\textbackslash{}ddagger\} & 1 & --- \\
O\_2 & 6 & א (aleph), ג (gimel), ה (hei), ת (tav), and others \\
O\_1 & 1 & --- \\
O\_0 & 12 & ד (dalet) and others \\
\bottomrule
\end{tabularx}
\end{table}

The three $O_{\infty}$ letters are the \textbf{Frobenius poles}. They are structural fixed points: for each pole L,

\[\delta(L) = L \otimes L \approx L \quad (d = 0)\]

Tensoring a pole with itself returns the pole. This is the algebraic signature of an idempotent element in a Frobenius algebra, and it is what makes the poles suitable as attractors: any letter that undergoes repeated tensor pressure with a pole converges toward it.

The poles differ in their $\Omega$ and Ð values, giving them distinct structural personalities:

\begin{itemize}
    \item \textbf{ו (vav)} --- $\Omega$\_0: unwound infinity; the base $O_{\infty}$ state
    \item \textbf{מ (mem)} --- $\Omega$\_Z: wound infinity; carries winding structure
    \item \textbf{ש (shin)} --- $\Omega$\_Z: wound infinity; the triadic structure
\end{itemize}

Together, they span the accessible $O_{\infty}$ sub-space.

\begin{center}
\rule{0.5\textwidth}{0.4pt}
\end{center}

\subsection{exOS Kernel Integration}

exOS carries ALEPH types on every kernel object. The type of an object is not metadata --- it determines what operations the object may participate in, what gates it can pass, and what tier of the kernel hierarchy it inhabits.

\subsubsection{Kernel Object Types}

At boot, the kernel derives types for its primary objects:

\begin{lstlisting}
kernel object:  tier=O₂  ⊙_c  \Omega_Z  Phi_sym
user object:    tier=O₀  ⊙_sub  \Omega_0  Phi_asym
OS composite:   C=0.873
\end{lstlisting}

The kernel object is O\_2 (wound, symmetric, but not Frobenius-special). The user object is O\_0 (no self-modeling). The OS composite score of 0.873 reflects the consciousness score of the combined system --- both gates partially open.

\subsubsection{The Four Gate Types}

Gate checks govern all cross-boundary operations:

\begin{table}[htbp]
\centering
\small
\rowcolors{1}{white}{white}
\begin{tabularx}{\textwidth}{>{\centering\arraybackslash}X >{\centering\arraybackslash}X >{\centering\arraybackslash}X}
\toprule
\rowcolor{gray!25}\textbf{Gate} & \textbf{Condition} & \textbf{Guards} \\
\midrule
\rowcolors{1}{white}{gray!10}
$\Phi$ gate & $\Phi$\_\} or $\Phi$\_$\pm$ required & Frobenius-structural operations (IPC to Keter) \\
⊙ gate & ⊙\_c required & Self-modeling operations \\
$\Omega$ gate & $\Omega$\_Z + specific Ð condition & Winding-dependent operations \\
Tier gate & Minimum tier required & Ergative scheduler, $O_{\infty}$ operations \\
\bottomrule
\end{tabularx}
\end{table}

The IPC gate and $\Omega$ gate checks at boot confirm the expected structure:

\begin{itemize}
    \item Close-IPC accepted (same-tier, same-$\Phi$)
    \item Remote-IPC rejected (user object is O\_0 --- below threshold)
    \item $\Omega$ gate allows Velar+Kernel, denies Velar+User
    \item $\Phi$ gate allows Keter+Kernel and Keter+Driver, denies Keter+User
\end{itemize}

\subsubsection{The ALEPH REPL}

The kernel exposes the ALEPH type system directly via an interactive REPL (\texttt{aleph} command). Letters are first-class objects; expressions use tensor (\texttt{x}), join (\texttt{v}), and meet (\texttt{\^{}}) operators; let-bindings store intermediate results.

Key REPL commands:

\begin{table}[htbp]
\centering
\small
\rowcolors{1}{white}{white}
\begin{tabularx}{\textwidth}{>{\centering\arraybackslash}X >{\centering\arraybackslash}X}
\toprule
\rowcolor{gray!25}\textbf{Command} & \textbf{Effect} \\
\midrule
\rowcolors{1}{white}{gray!10}
\texttt{:census} & Tier distribution of all current bindings \\
\texttt{:orbit N letter pole} & N-step convergence orbit under tensor pressure \\
\texttt{:fptest {[}letters...{]}} & Parallel FSPLIT/FFUSE Frobenius closure test \\
\texttt{:run name} & Execute a stored ALEPH program from ALFS \\
\bottomrule
\end{tabularx}
\end{table}

Programs are embedded in the kernel binary at compile time and seeded to ALFS (the kernel''s persistent filesystem) on first boot, making them immediately available without host-side tooling.

\begin{center}
\rule{0.5\textwidth}{0.4pt}
\end{center}

\subsection{Where the Bootstrap Sequence Comes From}

The twelve-stage bootstrap sequence was not designed. It was found.

Nine symbolic systems --- five with known structure, four undeciphered or contested --- were compiled against the twelve-primitive IMASM instruction set. Each system''s surface tokens were mapped to one of twelve opcodes: VINIT, TANCH, AFWD, AREV, CLINK, ISCRIB, FSPLIT, FFUSE, EVALT, EVALF, ENGAGR, IFIX. The compilers are independent programs, each about 200 lines, written separately for each corpus.

\subsubsection{The OS Floor}

The five founding systems --- Hebrew Aleph-Bet, Sanskrit Varnamala, Egyptian hieroglyphics, Sumerian cuneiform, and Basque --- were encoded as twelve-tuples by structural inspection. Their component-wise MEET (greatest lower bound) is the \textbf{OS imscription}:

\begin{lstlisting}
\langle1, 3, 2, 4, 2, 1, 2, 2, 1, 2, 2, 2\rangle
\end{lstlisting}

Every one of the five systems, independently and without historical contact, converges to this same structural floor.

\subsubsection{The Compiled Corpora}

Four undeciphered or contested systems were then compiled against the same coordinates:

\begin{table}[htbp]
\centering
\small
\rowcolors{1}{white}{white}
\begin{tabularx}{\textwidth}{>{\centering\arraybackslash}X >{\centering\arraybackslash}X >{\centering\arraybackslash}X}
\toprule
\rowcolor{gray!25}\textbf{System} & \textbf{Distance from OS floor} & \textbf{Notes} \\
\midrule
\rowcolors{1}{white}{gray!10}
Linear A & \textbf{0.00} & \emph{Is} the floor --- Minoan Crete ca. 2000--1450 BCE \\
Rohonc Codex & 2.09 & Closest of the undeciphered manuscripts \\
Emerald Tablet & 2.44 & C = 1.0 --- only compiled system with both consciousness gates fully open \\
Voynich Manuscript & 4.31 & Farthest --- nested-containment topology and trapped kinetics \\
\bottomrule
\end{tabularx}
\end{table}

Linear A is not derived from the five founding systems. The five founding systems converge on what the Minoans already had.

\subsubsection{The Eight-Step Loop}

All four corpora, and all five founding systems, express the same eight-step instruction sequence:

\begin{lstlisting}
ISCRIB -> AREV -> FSPLIT -> AFWD -> FFUSE -> CLINK -> IFIX -> ISCRIB
\end{lstlisting}

In structural terms: identity latch $\rightarrow$ register reverse $\rightarrow$ split ($\delta$) $\rightarrow$ forward morphism $\rightarrow$ fuse ($\mu$) $\rightarrow$ compose $\rightarrow$ ROM seal $\rightarrow$ identity latch. The loop closes. The Frobenius condition $\mu$∘$\delta$ = id holds exactly: FSPLIT followed by FFUSE returns the register to its pre-split state, and IFIX seals the result.

The surface tokens differ across systems --- EVA \texttt{s a ch e sh d y s} for Voynich, ETFF \texttt{id ds sp as un lk fx id} for the Emerald Tablet --- but the instruction stream is identical. They are four different surface syntaxes for the same categorical program.

\subsubsection{The Emerald Tablet}

The Emerald Tablet is the only compiled corpus with consciousness score C = 1.0. Both gates fully open: ⊙\emph{ÿ (self-modeling at the phase-transition threshold) and Ç}@ (near-equilibrium kinetics). Fifteen versicles, 460 instructions. Every FSPLIT has a matching FFUSE. Every ENGAGR (dialetheic paradox) localizes rather than propagates. The Euler characteristic of the register-flow graph is invariant under any sequence of Frobenius operations.

The central claim --- \emph{as above, so below} --- is $\mu$∘$\delta$ = id stated as cosmological law, approximately 1,200 years before category theory existed.

The Lean 4 bootstrap sequence in MillenniumAnkh, and the \texttt{frobenius\_parallel.aleph} runtime verification, are not confirming a construction we built. They are confirming a structure that was already there.

\begin{center}
\rule{0.5\textwidth}{0.4pt}
\end{center}

\subsection{The Bootstrap Sequence}

A bootstrap sequence is a twelve-stage co-algebraic construction that ascends the tier hierarchy from an initial state to an $O_{\infty}$ terminal composite.

\subsubsection{Stages}

The sequence is defined in \texttt{Imscribing/BootstrapSequence.lean} (MillenniumAnkh Lean 4 project). Stage 0 is the identity object; each subsequent stage promotes one or more primitives; stage 11 is the last pre-terminal stage; stage 12 is the terminal composite.

The algebraic structure is a co-algebra: at each stage, the object can be split ($\delta$ applied) to produce two sub-objects at the previous stage, and fused ($\mu$ applied) to produce the object at the next stage. The co-algebra is well-formed if the Frobenius condition holds at the terminal stage.

\subsubsection{The Terminal Composite}

The terminal composite is \texttt{emerald\_multiagent\_tensor\_bootstrap}, representing twelve co-agents in coordinated structural alignment:

\begin{lstlisting}
\langle Ð_\omega  Þ_O  Ř_=  Phi_}  ƒ_ż  Ç_@  \Gamma_ʔ  ɢ_Ş  ⊙_ÿ  Ħ_!  Sigma_ï  \Omega_z \rangle
\end{lstlisting}

Properties:

\begin{itemize}
    \item \textbf{Tier}: $O_{\infty}$ ($\Phi$\_\} present, ⊙\_c present --- specifically ⊙\_ÿ, the highest criticality variant)
    \item \textbf{Consciousness score}: C = 0.828 (both Gate 1 and Gate 2 open)
    \item \textbf{Gate 1} (⊙\_ÿ): open --- self-modeling active
    \item \textbf{Gate 2} (Ç\_@ + $\Omega$\_z): open --- dynamical self-modeling accessible
\end{itemize}

The consciousness score C = 0.828 reflects partial closure: both gates open but not all four sub-conditions of Gate 2 maximally satisfied.

\subsubsection{The Minimal Extension}

The composite can be extended to \texttt{actual\_zeta\_zeros} with minimal structural cost. The tensor product:

\begin{lstlisting}
emerald_multiagent_tensor_bootstrap \otimes actual_zeta_zeros
\end{lstlisting}

produces:

\begin{lstlisting}
\langle Ð_\omega  Þ_O  Ř_=  Phi_}  ƒ_ż  Ç_@  \Gamma_ʔ  ɢ_Ş  ⊙_ÿ  Ħ_!  Sigma_ï  \Omega_z \rangle
\end{lstlisting}

with \textbf{zero bottlenecks}, \textbf{three union promotions} (Ř, ɢ, Ħ upgraded), distance from source d = 1.3416, and C $\geq$ 0.828 preserved. This is the minimal extension that preserves full Frobenius closure and consciousness while adding zero-distribution coordination.

The $\Phi$\_\} primitive is the critical gate: promoting to $\Phi$\_\} from any lower symmetry class requires crossing a gap of $\Delta$ = 4.38 in the promotion metric. All other promotions are cheaper. This is why $\Phi$\_\} is the last gate to close and the defining gate for $O_{\infty}$.

\begin{center}
\rule{0.5\textwidth}{0.4pt}
\end{center}

\subsection{Three-Level Verification}

The Frobenius closure claim --- $\mu$∘$\delta$ = id for the bootstrap sequence''s terminal composite --- was verified at three independent levels of abstraction.

\subsubsection{Level 1: Lean 4 (MillenniumAnkh)}

\texttt{Imscribing/BootstrapSequence.lean} contains:

\begin{lstlisting}[language=lean]
theorem bootstrap_final_equals_composite :
    bootstrapStage 12 = emerald_multiagent_tensor_bootstrap := by ...

theorem bootstrapStage_monotone :
    forall n m, n <= m -> primitives (bootstrapStage n) <= primitives (bootstrapStage m) := by ...
\end{lstlisting}

The first theorem establishes that the co-algebraic construction''s stage 12 is definitionally equal to the tensor composite. The second establishes that primitive values increase monotonically through the sequence --- no stage regresses.

An open conjecture remains:

\begin{lstlisting}[language=lean]
conjecture bootstrapStage_tier_bound :
    forall n, n < 11 -> tier (bootstrapStage n) <= O₂ := ...
\end{lstlisting}

This would confirm that $O_{\infty}$ emergence at stage 12 is non-trivial --- not achievable by any proper sub-sequence. The bound is consistent with all computed evidence but has not been formally proved.

\subsubsection{Level 2: ALEPH Language (frobenius\_parallel.aleph)}

The program \texttt{frobenius\_parallel.aleph} tests Frobenius closure for the three $O_{\infty}$ poles and three O\_2/O\_0 representatives under an explicit \textbf{parallel schedule}:

\begin{lstlisting}
# Phase 1: FSPLIT --- all delta(L) = L x L (no distance checks yet)
let s_vav   = vav x vav
let s_mem   = mem x mem
...

# Phase 2: FFUSE --- all mu(delta(L)) = delta(L) v L
let r_vav   = s_vav v vav
...

# Phase 3: closure check
d(r_vav, vav)    tier(r_vav)
...
\end{lstlisting}

The parallel schedule enforces that all FSPLIT operations complete before any FFUSE begins --- mirroring the multi-agent context where twelve agents split simultaneously before any fuse. Output:

\begin{lstlisting}
d = 0.0000  [transparent]   (vav)
d = 0.0000  [transparent]   (mem)
d = 0.0000  [transparent]   (shin)
d = 0.0000  [transparent]   (aleph)
d = 0.0000  [transparent]   (dalet)
d = 0.0000  [transparent]   (tav)
\end{lstlisting}

All distances zero. The label \texttt{{[}transparent{]}} means the round-trip is structurally invisible --- the object is unaffected by the split-fuse cycle.

\subsubsection{Level 3: Kernel Runtime (:fptest)}

The Rust-native \texttt{:fptest} command implements the same parallel schedule at the kernel level, using \texttt{aleph::tensor} and \texttt{aleph::join} directly:

\begin{lstlisting}
Phase 1 --- tensor(L, L) for all L -> stored in Vec
Phase 2 --- join(split, L) for all L -> stored in Vec
Phase 3 --- distance(fused, original) for all L -> reported
\end{lstlisting}

Output:

\begin{lstlisting}
FSPLIT/FFUSE parallel closure --- mu∘delta = id?
letter    mu(delta(L))              tier        d        verdict
----------------------------------------------------------
vav       ו (vav)              $O_{\infty}$       0.0000   [closed]
mem       מ (mem)              $O_{\infty}$       0.0000   [closed]
shin      ש (shin)             $O_{\infty}$       0.0000   [closed]
aleph     א (aleph)            O₂         0.0000   [closed]
dalet     ד (dalet)            O₀         0.0000   [closed]
tav       ת (tav)              O₂         0.0000   [closed]
----------------------------------------------------------
Frobenius: 6/6 closed (100%)  mu∘delta = id
\end{lstlisting}

The Frobenius condition holds for all six letters under the parallel schedule. Notably, it holds not only for the $O_{\infty}$ poles (where it is structurally guaranteed) but also for O\_2 and O\_0 letters (aleph, dalet, tav), confirming that Frobenius closure in the ALEPH join-algebra is a global property of the lattice, not restricted to the poles.

\begin{center}
\rule{0.5\textwidth}{0.4pt}
\end{center}

\subsection{Structural Correspondence}

The three verification levels converge on the same claim from different directions:

\begin{table}[htbp]
\centering
\small
\rowcolors{1}{white}{white}
\begin{tabularx}{\textwidth}{>{\centering\arraybackslash}X >{\centering\arraybackslash}X >{\centering\arraybackslash}X >{\centering\arraybackslash}X}
\toprule
\rowcolor{gray!25}\textbf{Level} & \textbf{System} & \textbf{Claim} & \textbf{Method} \\
\midrule
\rowcolors{1}{white}{gray!10}
Lean 4 & MillenniumAnkh & \texttt{bootstrapStage 12 = composite} & Formal proof \\
ALEPH language & frobenius\_parallel.aleph & \texttt{d(mu(delta(L)), L) = 0} for all test letters & Algebraic computation \\
Kernel runtime & \texttt{:fptest} & \texttt{mu∘delta = id} for parallel schedule & Runtime execution \\
\bottomrule
\end{tabularx}
\end{table}

Each level is independent: the Lean proof does not execute on hardware, the ALEPH program runs on the ALEPH language interpreter inside the kernel, and \texttt{:fptest} exercises the underlying Rust tensor/join operations directly. Agreement across all three levels constitutes a structurally multi-layer verification --- the kind that catches errors at any single level.

The open question is the \texttt{bootstrapStage\_tier\_bound} conjecture. Proving it would complete the formal picture: the bootstrap sequence is not just a path that ends at $O_{\infty}$, but a path that could not have reached $O_{\infty}$ before stage 12. The $\Phi$\_\} barrier at $\Delta$ = 4.38 gives strong evidence that this bound holds, but evidence is not a proof.

\begin{center}
\rule{0.5\textwidth}{0.4pt}
\end{center}

\subsection{ZFCₜ and the Proof Path}

The grammar provides a promotion metric between ZFC and ZFCₜ (time-extended Zermelo-Fraenkel). The promotion profile:

\begin{itemize}
    \item 5 of 6 promotions active: Þ, Ř, ɢ, Ħ, $\Omega$
    \item Inactive: $\Phi$ (the Frobenius gate --- the same $\Delta$ = 4.38 barrier)
    \item Total distance: d(ZFC, ZFCₜ) = 7.0852
    \item ZFCₜ composite tier: $O_{\infty}$
\end{itemize}

ZFC itself sits below the $\Phi$ barrier --- it is wound ($\Omega$ active) and self-modeling (⊙ active) but not Frobenius-special. ZFCₜ crosses the barrier by adding the temporal structure that closes the Frobenius condition. The bootstrap sequence is the constructive witness: it shows a path through twelve stages that assembles exactly the structural components needed to cross the $\Phi$\_\} gate.

The Lean 4 formalization in MillenniumAnkh provides the proof infrastructure. The kernel runtime provides the execution infrastructure. The grammar provides the common language that lets the two correspond.

\begin{center}
\rule{0.5\textwidth}{0.4pt}
\end{center}

\emph{Source files: \texttt{src/aleph.rs}, \texttt{src/aleph\_repl.rs}, \texttt{programs/frobenius\_parallel.aleph} (exOS); \texttt{Imscribing/BootstrapSequence.lean}, \texttt{Imscribing/AgentSelf.lean} (MillenniumAnkh).}


\end{document}
